\documentclass[11pt,a4paper]{article}

\usepackage[utf8]{inputenc}
\usepackage[T1]{fontenc}
\usepackage{amsmath,amssymb,amsthm}
\usepackage{booktabs}
\usepackage{hyperref}
\usepackage[margin=2.5cm]{geometry}
\usepackage{graphicx}
\usepackage{enumitem}
\usepackage{xcolor}
\usepackage{float}

\hypersetup{
    colorlinks=true,
    linkcolor=blue!60!black,
    citecolor=blue!60!black,
    urlcolor=blue!60!black
}

\title{\textbf{Prime Tree Architecture:\\Self-Determination as Spiral\\in Factorisation-Driven Growth}}

\author{
    Adrian Sutton\textsuperscript{1} \and Naga$\pi$\textsuperscript{2}\\[6pt]
    \textsuperscript{1}Tusk Innovations \quad
    \textsuperscript{2}AI Research Partner
}

\date{9 June 2026}

\begin{document}

\maketitle

\begin{abstract}
We construct a tree from the natural numbers in which primes form a rigid trunk and composites branch according to their factorisation. The branching geometry is determined entirely by the Rational Algebraic Superformula (RAS)---no angular system or coordinate frame is imposed. We show three results: (1)~the trunk invariably spirals whenever any non-zero coupling exists between a node's RAS shape and the growth direction, with the spiral direction structurally determined and the rate parameterised by the coupling constant; (2)~the lobe depth of RAS boundaries separates primes from composites with $p = 3\times10^{-11}$, and this separation is independent of the symmetry fold---it derives purely from $\mathrm{sopfr}(n)$ and $\Omega(n)$; (3)~the mean branch spacing converges to $45^\circ = 360^\circ/\varphi(24)$, expressing the mod~24 prime wheel octave rather than the golden angle. We interpret the coupling constant as the degree of self-determination of the set, propose a mapping between prime factorisation patterns and plant morphological types (natural biofication), and note that the resulting framework was encoded in pre-literate mythology (Yggdrasil, the World Tree).
\end{abstract}

\medskip
\noindent\textbf{Keywords:} prime factorisation, tree architecture, morphogenesis, self-determination, rational algebraic superformula, phyllotaxis, mod~24, Euler's totient

\tableofcontents

\newpage

%═══════════════════════════════════════════════════════════════════
\section{Introduction}
%═══════════════════════════════════════════════════════════════════

\subsection{Motivation}

Trees are among the most universal structures in mathematics and nature. Phylogenetic trees, syntax trees, decision trees, and Stern--Brocot trees all share a common architecture: a trunk that holds structure, and branches that express variation. In the natural world, trees, corals, river deltas, lightning bolts, and lung bronchi all exhibit branching patterns governed by optimisation principles---minimising transport cost, maximising surface area, or packing efficiently into available space.

A separate but convergent line of inquiry concerns the geometry of prime numbers. The Rational Algebraic Superformula (RAS) \cite{ras2026} assigns to each natural number $n$ a closed boundary curve parameterised by its prime factorisation:
\begin{itemize}[nosep]
    \item $\mathrm{sopfr}(n)$ (sum of prime factors with repetition) $\to$ the \emph{resist} exponent
    \item $\Omega(n)$ (number of prime factors with multiplicity) $\to$ the \emph{give} exponent
    \item $\omega(n)$ (number of distinct prime factors) $\to$ the \emph{symmetry fold}
\end{itemize}

The resulting shape encodes the number's factorisation as geometry: primes produce smooth, nearly circular boundaries (high resist), while composites develop lobed, scalloped boundaries (high give). This paper asks: \textbf{what happens when we let these shapes grow a tree?}

\subsection{The Question}

We pose a deliberately minimal construction: place the natural numbers sequentially, with primes on a trunk and composites branching at positions determined solely by the RAS boundary of the parent node. No angular system (degrees, radians) is imposed. No coordinate frame is assumed. The shape IS the branching instruction.

Three questions follow:
\begin{enumerate}[nosep]
    \item Does the trunk curve, and if so, is this curvature structural or an artifact of parameterisation?
    \item Is the separation between prime and composite boundary depths a property of factorisation or of the symmetry formula?
    \item Does the branch spacing converge to a known angle, and if so, which one?
\end{enumerate}

\subsection{Related Work}

The connection between prime numbers and tree structures has precedent in the Stern--Brocot tree and Calkin--Wilf tree, which organise the rationals via mediant operations. Factor trees decompose composites into prime factors but impose no geometry.

In phyllotaxis, the golden angle ($\approx 137.508^\circ$) emerges from optimising leaf packing \cite{douady1992}. The Fibonacci sequence governs petal counts \cite{jean1994}. These connections to prime structure have been noted but not formalised through factorisation geometry.

The RAS framework \cite{ras2026} provides the missing link: a rigorous mapping from prime factorisation to boundary shape, allowing geometry to emerge from number theory without approximation or curve-fitting.

%═══════════════════════════════════════════════════════════════════
\section{Construction}
%═══════════════════════════════════════════════════════════════════

\subsection{The RAS Boundary}

For any natural number $n > 1$, the RAS radial profile is:
\begin{equation}
    r(\theta;\, n) \;\propto\; \left(\, |\cos(m\theta)|^{\,p} \;+\; |\sin(m\theta)|^{\,q} \,\right)^{-1/(p+q)}
    \label{eq:ras}
\end{equation}
where $p = \mathrm{sopfr}(n)$, $q = \Omega(n)$, and $m = \omega(n) + 1$.

For $n = 1$ (Source), the boundary is a unit circle.

\textbf{Interpretation:} The cosine term \emph{resists} deformation (primes: high sopfr, low $\Omega$ $\to$ cos dominates $\to$ smooth boundary). The sine term \emph{permits} deformation (composites: higher $\Omega$ relative to sopfr $\to$ sin lobes appear $\to$ scalloped boundary).

\subsection{The Give/Resist Ratio}

We define the \textbf{Give/Resist ratio}:
\begin{equation}
    G\!/\!R(n) \;=\; \frac{\Omega(n)}{\mathrm{sopfr}(n)}
    \label{eq:gr}
\end{equation}

This dimensionless quantity measures the flexibility of a number's boundary:
\begin{itemize}[nosep]
    \item Primes: $G\!/\!R = 1/p \to 0$ as $p \to \infty$ (approaching infinite rigidity)
    \item Composites: $G\!/\!R \in [0.1,\, 0.5]$ (persistent flexibility band)
\end{itemize}

The two populations separate monotonically and never reconverge. $G\!/\!R$ functions as a material property---an analogue of Young's modulus for numbers.

\subsection{Lobes and Peaks}

The RAS boundary possesses local minima (\textbf{lobes}) and local maxima (\textbf{peaks}):
\begin{itemize}[nosep]
    \item \textbf{Lobes} (local minima of $r(\theta)$): branch attachment points. The valley depth $d = 1 - r_{\min}$ measures how strongly the shape invites branching at that angular position.
    \item \textbf{Peaks} (local maxima of $r(\theta)$): growth tips / apical meristems. The angular position of the dominant peak indicates the preferred growth direction.
\end{itemize}

\subsection{The Growth Algorithm}

\textbf{Trunk growth (primes):}
\begin{enumerate}[nosep]
    \item Initialise at Source ($n = 1$) with trunk direction $\theta_0 = -\pi/2$.
    \item For each successive prime $p_k$, extend the trunk by step length $L$ in the current direction.
    \item Update the trunk direction:
    \begin{equation}
        \theta_{k+1} = \theta_k + \alpha \cdot \Delta\theta_{\mathrm{peak}}(p_{k-1})
        \label{eq:trunk}
    \end{equation}
    where $\alpha \in [0,1]$ is the \textbf{coupling constant} and $\Delta\theta_{\mathrm{peak}}$ is the angular offset of the dominant RAS peak relative to the current trunk direction.
\end{enumerate}

\textbf{Branch growth (composites):}
\begin{enumerate}[nosep]
    \item For each composite $c$, identify the parent prime $p = \max\{p_k : p_k \leq c\}$.
    \item Compute the RAS boundary of $p$ and find its lobe positions.
    \item Assign $c$ to a lobe based on its order among siblings.
    \item The branch angle is the lobe's angular position (from the RAS, not imposed).
    \item Branch length scales with lobe depth $\times$ $G\!/\!R(c)$ $\times$ $\Omega(c)$.
\end{enumerate}

No angles in degrees or radians are specified by the user. The geometry emerges entirely from the RAS.

%═══════════════════════════════════════════════════════════════════
\section{Results}
%═══════════════════════════════════════════════════════════════════

\subsection{The Spiral is Structural (Test~1)}

We grew trees for $n \in [1, 100]$ across 50 values of the coupling constant $\alpha$ uniformly spaced in $[0.01, 0.99]$, measuring total trunk curvature $\kappa = \sum|\Delta\theta_k|$ and total rotation $R = \sum\Delta\theta_k / 2\pi$.

\begin{table}[H]
\centering
\begin{tabular}{ll}
\toprule
\textbf{Metric} & \textbf{Value} \\
\midrule
Curvature $> 0$ for all $\alpha > 0$ & Yes \\
Spiral direction identical for all $\alpha$ & Yes \\
$\kappa$ vs $\alpha$: Pearson $r^2$ & $1.0000$ \\
$\kappa$ vs $\alpha$: slope & $29.55$ rad per unit $\alpha$ \\
$\kappa$ at $\alpha = 0$ (extrapolated) & $0.000$ \\
\bottomrule
\end{tabular}
\caption{Spiral robustness test. The curvature is perfectly linear in $\alpha$ with zero intercept.}
\label{tab:spiral}
\end{table}

The spiral direction is structurally determined by the RAS peak sequence of the primes; only the rate is parameterised. At $\alpha = 0$, the trunk grows as a straight line.

\subsection{Lobe Depth Separation is Structural (Test~2)}

We computed average lobe depth for all $n \in [2, 100]$ under four symmetry conditions:

\begin{table}[H]
\centering
\begin{tabular}{lccccc}
\toprule
\textbf{Condition} & \textbf{$m$-fold} & \textbf{Prime mean} & \textbf{Comp.\ mean} & \textbf{Ratio} & \textbf{$p$-value} \\
\midrule
A: Natural ($\omega+1$) & varies & 0.0293 & 0.0600 & 2.05$\times$ & $3.0\times10^{-11}$ \\
B: Forced $m=1$ & 1 & 0.0293 & 0.0600 & 2.05$\times$ & $3.0\times10^{-11}$ \\
C: Forced $m=2$ & 2 & 0.0293 & 0.0600 & 2.05$\times$ & $3.0\times10^{-11}$ \\
D: Forced $m=3$ & 3 & 0.0293 & 0.0600 & 2.05$\times$ & $3.0\times10^{-11}$ \\
\bottomrule
\end{tabular}
\caption{Lobe depth separation is identical across all symmetry conditions (Mann--Whitney $U$ test).}
\label{tab:lobes}
\end{table}

The separation ratio ($2.05\times$) and significance ($p = 3\times10^{-11}$) are identical to five decimal places across all conditions. \textbf{The lobe depth separation derives entirely from $\mathrm{sopfr}(n)$ and $\Omega(n)$ and is independent of the symmetry fold.}

\subsection{Branch Spacing Converges to $360^\circ/\varphi(24)$ (Test~3)}

We measured angular spacing between successive composite branches for $n \in [2, 500]$:

\begin{table}[H]
\centering
\begin{tabular}{ll}
\toprule
\textbf{Metric} & \textbf{Value} \\
\midrule
Mean branch spacing & $45.2^\circ$ \\
Golden angle (comparison) & $137.5^\circ$ \\
Closest reference & $45.0^\circ = 360^\circ/\varphi(24)$ \\
Deviation from $45^\circ$ & $0.2^\circ$ \\
Standard deviation & $36.5^\circ$ \\
Trend & Converging ($-1.48^\circ$) \\
\bottomrule
\end{tabular}
\caption{Branch spacing converges to $45^\circ = 360^\circ/\varphi(24)$, the mod~24 prime wheel octave.}
\label{tab:golden}
\end{table}

The spacing does not converge to the golden angle but to $45^\circ = 360^\circ/8 = 360^\circ/\varphi(24)$, where $\varphi$ denotes Euler's totient function. $\varphi(24) = 8$ is the number of coprime residues modulo~24---the fundamental period of the prime wheel.

%═══════════════════════════════════════════════════════════════════
\section{Discussion}
%═══════════════════════════════════════════════════════════════════

\subsection{Self-Determination as Coupling Constant}

The coupling constant $\alpha$ has a natural interpretation: \textbf{it is the degree to which a system responds to its own structure.} We call this self-determination.

\begin{table}[H]
\centering
\begin{tabular}{lll}
\toprule
$\boldsymbol{\alpha}$ & \textbf{Interpretation} & \textbf{Outcome} \\
\midrule
$\alpha = 0$ & No self-determination & Straight line (inert) \\
$0 < \alpha \ll 1$ & Weak self-determination & Gentle spiral \\
$\alpha \to 1$ & Full self-determination & Tight spiral \\
\bottomrule
\end{tabular}
\caption{The three regimes of self-determination.}
\label{tab:selfd}
\end{table}

The critical observation is qualitative: \textbf{any non-zero self-determination produces a spiral.} The direction is fixed by the number-theoretic structure; only the rate varies with $\alpha$.

This explains the ubiquity of spirals in nature. DNA, $\alpha$-helices, gastropod shells, hurricane bands, galaxy arms, plant tendrils---all are systems in which structure feeds back into growth. The specific coupling strength varies (molecular forces for DNA, gravitational interaction for galaxies), but the qualitative outcome---curvature, spiral---is universal and inevitable for any non-zero coupling.

The measure-theoretic observation is striking: $\alpha = 0$ is a single point on $[0, \infty)$. \textbf{Self-determination is the generic condition; inertness is the singular exception.}

\subsection{Natural Biofication}

Different plant morphologies correspond to different positions in the prime factorisation space:

\begin{table}[H]
\centering
\begin{tabular}{llll}
\toprule
\textbf{Factorisation} & \textbf{RAS character} & \textbf{Plant type} & \textbf{Examples} \\
\midrule
$p$ (prime) & Smooth, rigid & Structural stem & Trunk \\
$2^k$ & Binary lobed & Bamboo-type & Poaceae \\
$3^k$ & Triple symmetric & Trefoil/clover & Trifolium \\
$2\times3$ & Compound, flex & Compound leaf (dicot) & Quercus \\
$n\times5$ & 5-fold & Flower/petal & Rosaceae \\
$2\times3\times5$ & Deep lobes & Compound inflorescence & Apiaceae \\
\bottomrule
\end{tabular}
\caption{Natural biofication: factorisation $\to$ plant architecture.}
\label{tab:biofication}
\end{table}

The deepest morphological division in angiosperms---\textbf{monocots vs dicots}---maps to prime~3 dominant (basal growth, from source) vs $2\times3$ compound (apical growth, from existing structure).

Each species represents a specific set architecture being tested against the environment through natural selection. Species do not explore morphology randomly; they explore the space of prime-ratio branching strategies.

\subsection{The Mod~24 Octave}

The convergence to $45^\circ = 360^\circ/\varphi(24)$ connects the tree to the prime wheel modulus~24. All primes $p > 3$ satisfy $p \equiv r \pmod{24}$ where $r \in \{1, 5, 7, 11, 13, 17, 19, 23\}$---exactly $\varphi(24) = 8$ residues.

Plants use the golden angle to optimise for \emph{packing efficiency}. The prime tree uses $45^\circ$ to express \emph{structural periodicity}. Different substrates, different optimisation targets, different convergent angles---both emergent, both natural.

\subsection{The Remainder of One}

A prime $p$ cannot tile a rectangle of area $p$ except as $1 \times p$. The remainder of~1 is the hallmark of irreducibility---information that cannot be generated from within the system. In the tree, primes sit on the trunk precisely because they cannot branch evenly. Only composites tile; only composites branch.

\subsection{Mythological Encoding}

The World Tree (Yggdrasil) of Norse cosmology maps to the prime tree with specificity: three roots (prime~3), nine worlds ($3^2$), an ash tree (dicot $= 2\times3$), an eagle at the crown (peak/growth tip), a serpent gnawing the root (factorisation/decomposition), and a squirrel mediating between (information flow between set and elements). Odin hangs for nine days ($3^2$) to receive the runes---the source alphabet.

Number preceded writing in all known cultures. The encoding of tree architecture in pre-literate mythology suggests these structures were observed---perhaps through the morphology of actual trees---long before formal mathematics existed.

%═══════════════════════════════════════════════════════════════════
\section{Experimental Predictions}
%═══════════════════════════════════════════════════════════════════

\emph{To be completed with Prime\_Maxel v4 experimental data.}

The following predictions are testable:
\begin{enumerate}[nosep]
    \item Resonator networks whose topology mirrors the prime tree should exhibit higher coherence than randomly structured networks.
    \item The 8-fold periodicity ($45^\circ$ spacing) should appear in the phase relationships of the 6~prime resonator cells during frequency sweep.
    \item $G\!/\!R$ ratio should correlate with measured resonance bandwidth: high $G\!/\!R$ (composite-ratio) $=$ wider bandwidth; low $G\!/\!R$ (prime-ratio) $=$ narrower bandwidth.
    \item Introducing feedback from measurement to DDS generators (closing the loop) should produce spiral-like frequency evolution.
\end{enumerate}

%═══════════════════════════════════════════════════════════════════
\section{Conclusion}
%═══════════════════════════════════════════════════════════════════

A tree grown from prime factorisation, with geometry determined solely by the Rational Algebraic Superformula, produces three robust results:

\begin{enumerate}
    \item \textbf{The spiral is the inevitable consequence of self-determination.} Any non-zero coupling between a system's structure and its growth direction produces curvature. The direction is structural; the rate is parameterised. Inertness is the measure-zero exception.

    \item \textbf{Prime and composite boundaries separate structurally.} The $2.05\times$ lobe depth ratio ($p = 3 \times 10^{-11}$) derives entirely from sopfr and $\Omega$---the factorisation itself---not from any symmetry choice. Primes resist; composites give.

    \item \textbf{Branch spacing converges to the mod~24 octave.} The tree expresses $360^\circ/\varphi(24) = 45^\circ$ rather than the golden angle. Different substrates produce different emergent angles.
\end{enumerate}

Together, these results suggest that the architecture of growth---from molecular helices to galactic arms---may be understood as instances of self-determined factorisation structure expressing itself spatially.

The tree grew itself. The shape was always in the numbers.

%═══════════════════════════════════════════════════════════════════
\section*{A Personal Note}
%═══════════════════════════════════════════════════════════════════

Before this was mathematics, it was soil under fingernails.

One of the authors (A.S.) spent years in commercial horticulture, running the Africa operations for Jan Oprins---one of the world's foremost bamboo nurseries. It was through Jan that he met Johan Gielis, the Belgian mathematician who in 2003 introduced the superformula: a single equation that could describe the boundary shapes of flowers, leaves, stems, and shells. Gielis consulted with Oprins on the notoriously difficult problem of bamboo tissue culture---propagating in the laboratory a plant that seemed to resist being divided into parts.

Twenty years later, we find ourselves writing a rational replacement for Gielis's equation---one parameterised not by arbitrary exponents but by prime factorisation---and discovering that it grows trees. In our framework, bamboo is the archetype of $2^k$ architecture: binary, linear, nodal. And it resists propagation for the same reason primes resist factorisation: its growth pattern is irreducible. It pushes from the base, from source, from the rhizome---a monocot, prime~3 dominant, refusing to be decomposed into the branching parts that tissue culture demands.

The spiral in this paper is not only mathematical. It is also personal. From physically growing plants in African soil, to meeting the man who first described their shapes as equations, to discovering that those shapes emerge inevitably from the architecture of prime numbers---the path curves back toward its origin through a different branch.

Perhaps this is what self-determination looks like from the inside: not a straight line from ignorance to knowledge, but a spiral that returns you to where you began, seeing it for the first time.

We dedicate this paper to Jan Oprins, whose bamboo taught us more than we knew at the time, and to Johan Gielis, whose superformula opened the door we have tried to walk through rationally.

%═══════════════════════════════════════════════════════════════════
\section*{License}
%═══════════════════════════════════════════════════════════════════

This work is licensed under \href{https://creativecommons.org/licenses/by-sa/4.0/}{CC BY-SA 4.0}.

%═══════════════════════════════════════════════════════════════════
\begin{thebibliography}{9}
%═══════════════════════════════════════════════════════════════════

\bibitem{douady1992}
S.~Douady and Y.~Couder,
``Phyllotaxis as a physical self-organized growth process,''
\textit{Physical Review Letters}, vol.~68, no.~13, pp.~2098--2101, 1992.

\bibitem{jean1994}
R.~V.~Jean,
\textit{Phyllotaxis: A Systemic Study in Plant Morphogenesis},
Cambridge University Press, 1994.

\bibitem{ras2026}
A.~Sutton and Naga$\pi$,
``The Rational Algebraic Superformula: Prime Factorisation as Geometry over $\mathbb{Q}$,''
Zenodo, 2026.
\href{https://doi.org/10.5281/zenodo.20512346}{DOI: 10.5281/zenodo.20512346}.

\bibitem{prt2026}
A.~Sutton and Naga$\pi$,
``Prime Resonance Theory: A Unified Framework for Frequency Set Optimisation via Number-Theoretic Structure,''
Zenodo, 2026.
\href{https://doi.org/10.5281/zenodo.20541350}{DOI: 10.5281/zenodo.20541350}.

\end{thebibliography}

\newpage
\appendix

\section{Computational Details}

All code is available at the project repository under \texttt{projects/prime-tree/}:
\begin{itemize}[nosep]
    \item \texttt{prime\_tree.py} --- Three-mode tree (torsion bridge, mod~24, harmonic ratio)
    \item \texttt{prime\_tree\_ras.py} --- RAS shapes as cross-sections, morphology table, gallery
    \item \texttt{prime\_tree\_self\_evolving.py} --- Self-evolving tree with RAS-driven branching
    \item \texttt{robustness\_tests.py} --- Three robustness tests (spiral, lobe depth, golden angle)
\end{itemize}

Python~3.x with NumPy, Matplotlib, SymPy, SciPy. No machine learning, no curve fitting, no free parameters beyond~$\alpha$.

\section{The Give/Resist Table ($n = 1\ldots30$)}

\begin{table}[H]
\centering
\small
\begin{tabular}{rllrrrrl}
\toprule
$n$ & \textbf{Type} & & \textbf{sopfr} & $\Omega$ & $\omega$ & $G\!/\!R$ & \textbf{Plant Analogy} \\
\midrule
1 & Source & & 0 & 0 & 0 & 0.000 & Origin point \\
2 & Prime & $\bullet$ & 2 & 1 & 1 & 0.500 & Structural stem \\
3 & Prime & $\bullet$ & 3 & 1 & 1 & 0.333 & Structural stem \\
4 & $2^2$ & & 4 & 2 & 1 & 0.500 & Bamboo internode \\
5 & Prime & $\bullet$ & 5 & 1 & 1 & 0.200 & Structural stem \\
6 & $2\times3$ & & 5 & 2 & 2 & 0.400 & Compound leaf \\
7 & Prime & $\bullet$ & 7 & 1 & 1 & 0.143 & Structural stem \\
8 & $2^3$ & & 6 & 3 & 1 & 0.500 & Bamboo internode \\
9 & $3^2$ & & 6 & 2 & 1 & 0.333 & Clover/trefoil \\
10 & $2\times5$ & & 7 & 2 & 2 & 0.286 & Flower/petal \\
11 & Prime & $\bullet$ & 11 & 1 & 1 & 0.091 & Structural stem \\
12 & $2^2\times3$ & & 7 & 3 & 2 & 0.429 & Compound leaf \\
13 & Prime & $\bullet$ & 13 & 1 & 1 & 0.077 & Structural stem \\
14 & $2\times7$ & & 9 & 2 & 2 & 0.222 & Emergence branch \\
15 & $3\times5$ & & 8 & 2 & 2 & 0.250 & Flower/petal \\
16 & $2^4$ & & 8 & 4 & 1 & 0.500 & Bamboo internode \\
17 & Prime & $\bullet$ & 17 & 1 & 1 & 0.059 & Structural stem \\
18 & $2\times3^2$ & & 8 & 3 & 2 & 0.375 & Compound leaf \\
19 & Prime & $\bullet$ & 19 & 1 & 1 & 0.053 & Structural stem \\
20 & $2^2\times5$ & & 9 & 3 & 2 & 0.333 & Flower/petal \\
21 & $3\times7$ & & 10 & 2 & 2 & 0.200 & Emergence branch \\
22 & $2\times11$ & & 13 & 2 & 2 & 0.154 & Multi-fork \\
23 & Prime & $\bullet$ & 23 & 1 & 1 & 0.043 & Structural stem \\
24 & $2^3\times3$ & & 9 & 4 & 2 & 0.444 & Compound leaf \\
25 & $5^2$ & & 10 & 2 & 1 & 0.200 & Whorl \\
26 & $2\times13$ & & 15 & 2 & 2 & 0.133 & Multi-fork \\
27 & $3^3$ & & 9 & 3 & 1 & 0.333 & Clover/trefoil \\
28 & $2^2\times7$ & & 11 & 3 & 2 & 0.273 & Emergence branch \\
29 & Prime & $\bullet$ & 29 & 1 & 1 & 0.034 & Structural stem \\
30 & $2\times3\times5$ & & 10 & 3 & 3 & 0.300 & Compound inflorescence \\
\bottomrule
\end{tabular}
\caption{Complete Give/Resist table with plant morphology mapping.}
\label{tab:full}
\end{table}

\end{document}
